
\section{引言}

从视觉输入中感知和理解三维空间信息是人类空间智能的基石~\citep{arterberry2000perception}，也是机器人技术和混合现实等应用的关键需求。这一基本能力催生了广泛的3D视觉任务，包括单目深度估计~\citep{eigen2014depth}、运动恢复结构~\citep{snavely2006photo}、多视图立体视觉~\citep{seitz2006comparison}和同步定位与建图~\citep{mur2015orb}。尽管这些任务在概念上具有很强的重叠性——通常仅相差单一因素，例如输入视图的数量——但主流范式一直是为每个任务开发高度专门化的模型。虽然近期的努力~\citep{wang2024dust3r,wang2025vggt}已探索通过统一模型来同时处理多个任务，但它们通常存在关键局限性：往往依赖复杂、特制的架构，从头开始通过任务联合优化进行训练，因而无法有效利用大规模预训练模型。

在本工作中，我们跳出既有的3D任务定义，回归到受人类空间智能启发的更基本目标：从任意视觉输入中恢复3D结构，无论是单张图像、场景的多个视图还是视频流。我们摒弃复杂的架构工程，以两个核心问题为导向，追求一种极简的建模策略。第一，\textit{是否存在\textbf{最小的预测目标集合}，还是需要对众多3D任务进行联合建模？}第二，\textit{\textbf{单一的朴素Transformer}}是否足以实现这一目标？我们的工作对这两个问题都给出了肯定答案。我们提出\paperfullname，这是一个单一Transformer模型，通过一种特殊选择的射线表示，专门针对联合\textbf{任意视角深度与位姿估计}进行训练。我们证明，这种极简方法足以从任意数量的图像中重建视觉空间，无论相机位姿是否已知。

\paperfullname 将上述几何重建目标表述为一个密集预测任务。对于给定的 $N$ 张输入图像，模型经过训练输出 $N$ 张对应的深度图和射线图，每张图与其各自输入像素对齐。实现这一目标的架构始于一个标准的预训练视觉Transformer（例如，~\citealt{oquab2023dinov2}）作为骨干网络，利用其强大的特征提取能力。为了处理任意数量的视图，我们引入了一个关键改进：一种输入自适应的跨视角自注意力机制。该模块在选定层的前向传播过程中动态重排标记，从而实现所有视图间的高效信息交换。对于最终预测，我们提出了一个新的双DPT头，通过以不同的融合参数处理同一组特征，同时输出深度值和射线值。为了增强灵活性，模型可以通过一个简单的相机编码器可选地融入已知的相机位姿，使其能够适应各种实际场景。这一整体设计产生了一个简洁且可扩展的架构，直接继承了其预训练骨干网络的缩放特性。

我们通过教师-学生训练范式来训练Depth Anything 3，以统一多样化的训练数据，这对于通用模型是必要的。我们的数据来源包括多种格式，如真实世界的深度相机采集数据（例如，\citealt{baruch2021arkitscenes}）、3D重建数据（例如，\citealt{reizenstein2021common}）以及合成数据，其中真实世界的深度可能质量较差（\cref{fig:poordataset}）。为了解决这个问题，我们采用了受先前工作~\citep{yang2024depth,yang2024depthv2}启发的伪标签策略。具体来说，我们在合成数据上训练一个强大的教师单目深度模型，为所有真实世界数据生成密集、高质量的伪深度。至关重要的是，为了保持几何完整性，我们将这些密集的伪深度图与原始的稀疏或噪声深度进行对齐。这种方法被证明非常有效，在保证几何精度的同时，显著增强了标签的细节和完整性。

为了更好地评估我们的模型并跟踪该领域的进展，我们建立了一个全面的基准，用于评估几何和位姿精度。该基准包含5个不同的数据集，总计超过89个场景，涵盖从物体级别到室内和室外环境。通过跨场景直接评估位姿精度，并将预测的位姿和深度融合成3D点云进行精度评估，该基准忠实地衡量了视觉几何估计器的位姿和深度精度。实验表明，我们的模型在20个设置中有18个达到了最先进水平的性能。此外，在标准的单目基准上，我们的模型优于Depth Anything 2~\citep{yang2024depthv2}。

为了进一步展示\longname在推动其他3D视觉任务方面的基本能力，我们引入了一个具有挑战性的前馈新视角合成（FF-NVS）基准，包含超过160个场景。我们遵循极简建模策略，用额外的DPT头微调我们的模型，以预测像素对齐的3D高斯参数。大量实验产生了两个关键发现：1）将几何基础模型微调用于NVS，显著优于高度专门化的任务特定模型~\citep{xu2025depthsplat}；2）增强的几何重建能力与改进的FF-NVS性能直接相关，确立了\longname作为该任务最优骨干网络的地位。
